\section{第2課 これは　本です}

\subsection{基本課文}

\begin{enumerate}
    \item これは　本です。
    \item それは　何です。
    \item あれは　だれの　傘ですか。
    \item この　カメラは　スミスさんのです。
\end{enumerate}

\begin{itemize}
    \item 甲：これは　テレビですか。
    \item 乙：いいえ、それは　テレビではありません。　パソコンです。
\end{itemize}

\begin{itemize}
    \item 甲：それは　何ですか。
    \item 乙：これは　日本語の　本です。
\end{itemize}

\begin{itemize}
    \item 甲：森さんの　かばんは　どれですか。
    \item 乙：あの　かばんです。
\end{itemize}

\begin{itemize}
    \item 甲：その　ノートは　誰のですか。
    \item 乙：私のです。
\end{itemize}


\subsection{语法}

\begin{itemize}
    \item これ、それ、あれ　は　〜　です。
    \item だれですか、どなたですか、何ですか。
    \item どれ、どの
\end{itemize}


\subsection{表达}

\begin{itemize}
    \item 方
    \item 何歳ですか、おいくつですか。
    \item どうぞ。
    \item どうも、ありがとう　ございます。
    \item 100以下数字
    \item 家族称谓
\end{itemize}

\subsection{应用课文:家族の写真}

小野：李さん、これは　何ですか。

李：これですか。家族の写真です。

小野：この　方は　どなたですか。

李：わたしの　毋です。

小野：お母さんは　おいくつですか。

李：５２歳です。

\vspace{1em}

李：小野さん、これ、どうぞ。

小野：えっ、何ですか。

李：お土産です。

小野：わあ、シルクの　ハンカチですか。

李：ええ、汕頭の　ハンカチです。　中国の　名産品です。

小野：どうも　ありがとう　ございます。


\begin{vocablist}[2]
    \vocab{本}
    \vocab{かばん}
    \vocab{ノート}
    \vocab{鉛筆}
    \vocab{傘}
    \vocab{靴}
    \vocab{新聞}
    \vocab{雑誌}
    \vocab{辞書}
    \vocab{カメラ}
    \vocab{テレビ}
    \vocab{パソコン}
    \vocab{ラジオ}
    \vocab{電話}
    \vocab{机}
    \vocab{椅子}
    \vocab{鍵}
    \vocab{時計}
    \vocab{手帳}
    \vocab{写真}
    \vocab{車}
    \vocab{自転車}
    \vocab{お土産}
    \vocab{名産品}
    \vocab{シルク}
    \vocab{スカーフ}
    \vocab{ハンカチ}
    \vocab{会社}
    \vocab{方}
    \vocab{人}
    \vocab{家族}
    \vocab{母}
    \vocab{お母さん}
    \vocab{日本語}
    \vocab{中国語}
    \vocab{これ、それ、あれ、どれ}
    \vocab{何、誰、どなた}
    \vocab{この、その、あの、どの}
    \vocab{えっ}
    \vocab{わあ}
    \vocab{ええ}
    \vocab{長嶋}
    \vocab{日本}
    \vocab{汕頭}
    \vocab{ロンドン}
    \vocab{おいくつ}
    \vocab{何歳}
\end{vocablist}
